\section{Experimental Procedure}\label{sec:task2-procedure}

\begin{enumerate}
    \item Prepare the planar checkerboard, measure the physical square size, and record the number of inner corners in each direction.
    \item Capture clear images at different distances, translations, and tilt angles. Record the phone model, rear-camera use, image resolution, and acquisition settings.
    \item Read each image and check whether it can be loaded. A read failure is recorded for that file and does not stop processing of the remaining images.
    \item Convert a color image to grayscale with \texttt{cv2.cvtColor}, then detect checkerboard inner corners with \texttt{cv2.findChessboardCorners}. Record the two-dimensional corners for successful images and the filename and reason for failures.
    \item Refine detected corners to subpixel accuracy with \texttt{cv2.cornerSubPix}.
    \item Generate planar three-dimensional object points from the checkerboard configuration and square size, and pair them with the two-dimensional image points. The geometry follows Equation~\ref{eq:task2-homography}.
    \item Check the number of valid images and the consistency of image sizes. Mixed resolutions and fewer valid views than \texttt{--min-images} must be rejected.
    \item Call \texttt{cv2.calibrateCamera} to estimate the camera matrix, distortion coefficients, and the rotation and translation vectors for each valid view. The camera matrix is defined in Equation~\ref{eq:task2-intrinsic}; the distortion model is given in Equations~\ref{eq:task2-radial-distortion} and~\ref{eq:task2-tangential-distortion}.
    \item Reproject the three-dimensional points with \texttt{cv2.projectPoints}, compute pointwise errors using Equation~\ref{eq:task2-point-reprojection}, and summarize the three error definitions in Section~\ref{sec:task2-results}.
    \item Save the camera matrix, distortion coefficients, poses, error statistics, image counts, and detection status to the NPZ result. If requested, use \texttt{cv2.undistort} to save corrected images.
    \item Use the recorded result and the criteria in Section~\ref{sec:task2-analysis} to assess parameter consistency and data quality.
\end{enumerate}

The complete control flow implemented by the OpenCV program is summarized in Figure~\ref{fig:calibration-workflow}.

\begin{figure}[htbp]
    \centering
    \includegraphics[width=0.96\linewidth]{calibration_workflow.pdf}
    \caption{OpenCV camera-calibration workflow implemented in this experiment.}
    \label{fig:calibration-workflow}
\end{figure}
